\documentclass{article}

% ============================================================
% A1 POSTER - CANONICAL REPOSITORY SOURCE
% Compile from poster/ after generating result macros and figures.
% ============================================================

\usepackage[utf8]{inputenc}
\usepackage[T1]{fontenc}
\usepackage[
  paperwidth=594mm,
  paperheight=841mm,
  top=16mm,
  bottom=16mm,
  left=18mm,
  right=18mm
]{geometry}

\usepackage{helvet}
\renewcommand{\familydefault}{\sfdefault}
\usepackage{xcolor}
\usepackage{graphicx}
\usepackage{ragged2e}
\usepackage{microtype}
\usepackage{multicol}
\usepackage{needspace}
\usepackage{tikz}
\usetikzlibrary{arrows.meta,positioning,calc}

\pagestyle{empty}
\setlength{\parindent}{0pt}
\setlength{\parskip}{5.5mm}
\setlength{\columnsep}{10mm}
\hyphenpenalty=10000
\exhyphenpenalty=10000

% ---------- Colour palette ----------
\definecolor{ink}{HTML}{183044}
\definecolor{trierblue}{HTML}{0077B5}
\definecolor{deepblue}{HTML}{005A87}
\definecolor{accent}{HTML}{D36B2C}
\definecolor{lightblue}{HTML}{EAF5FA}
\definecolor{lightgray}{HTML}{F2F5F7}
\definecolor{midgray}{HTML}{CBD5DC}
\definecolor{darkgray}{HTML}{526575}
\color{ink}

% ---------- Typography ----------
\newcommand{\TargetAccuracy}{50.0\%}
\newcommand{\ContextAccuracy}{55.5\%}
\newcommand{\BertAccuracy}{67.1\%}
\newcommand{\TargetAccuracyFraction}{0.500000}
\newcommand{\ContextAccuracyFraction}{0.554859}
\newcommand{\BertAccuracyFraction}{0.670846}
\newcommand{\TargetBarValue}{50.0}
\newcommand{\ContextBarValue}{55.5}
\newcommand{\BertBarValue}{67.1}
\newcommand{\TargetCorrect}{319}
\newcommand{\BertCorrect}{428}
\newcommand{\ContextCorrect}{354}
\newcommand{\GainLower}{6.6}
\newcommand{\GainUpper}{16.9}
\newcommand{\AccuracyGain}{11.6}
\newcommand{\BertMacroF}{66.7\%}
\newcommand{\ContextMacroF}{55.1\%}
\newcommand{\BertAuc}{0.716}
\newcommand{\DifferenceLower}{-16.9}
\newcommand{\DifferenceUpper}{-6.6}
\newcommand{\MacroDifferenceLower}{-16.9}
\newcommand{\MacroDifferenceUpper}{-6.6}
\newcommand{\BertGainCorrect}{74}
\newcommand{\BertSameRecall}{77.4\%}
\newcommand{\BertDifferentRecall}{56.7\%}
\newcommand{\ContextSameRecall}{46.4\%}
\newcommand{\ContextDifferentRecall}{64.6\%}
\newcommand{\TrainCount}{5,428}
\newcommand{\ValidationCount}{638}
\newcommand{\TestCount}{1,400}
\newcommand{\TuneCount}{4,342}
\newcommand{\HoldoutCount}{1,086}
\newcommand{\ExperimentSeed}{2026}
\newcommand{\GloveDimensions}{300}
\newcommand{\ContextCandidates}{$\pm 2$, $\pm 5$, $\pm 10$, sentence}
\newcommand{\SelectedContextWindow}{$\pm 2$}
\newcommand{\BootstrapResamples}{1,000}
\newcommand{\BestAccuracyLayer}{9}
\newcommand{\FinalLayerNumber}{12}
\newcommand{\BertOnlyCount}{178}
\newcommand{\BothCorrectCount}{250}
\newcommand{\StaticOnlyCount}{104}
\newcommand{\BothWrongCount}{106}
\newcommand{\BertErrorCount}{210}

\input{authors.tex}
\newcommand{\bodyfont}{\fontsize{24}{27}\selectfont}
\newcommand{\smallfont}{\fontsize{22}{26}\selectfont}
\newcommand{\tinyposterfont}{\fontsize{21}{25}\selectfont}
\newcommand{\sectiontitle}[1]{%
  \Needspace{4\baselineskip}%
  \par\vspace{5mm}%
  {\fontsize{30}{35}\selectfont\bfseries\textcolor{trierblue}{#1}\par}%
  \vspace{6mm}%
}
\newcommand{\subsectiontitle}[1]{%
  \Needspace{3\baselineskip}%
  \par\vspace{4mm}%
  {\fontsize{25}{29}\selectfont\bfseries\textcolor{deepblue}{#1}\par}%
  \vspace{4mm}%
}
\newcommand{\thinrule}{\par\vspace{2mm}{\color{midgray}\rule{\linewidth}{0.55pt}}\par\vspace{2mm}}

% ---------- Small highlight box ----------
\newcommand{\keybox}[1]{%
  \colorbox{lightblue}{%
    \parbox{0.955\linewidth}{\vspace{2mm}\RaggedRight #1\vspace{2mm}}%
  }%
}

% ---------- Result bar helper ----------
% #1 = label, #2 = percentage width on a true 0-100 scale, #3 = printed result, #4 = fill colour
\newcommand{\resultbar}[4]{%
  \begin{tikzpicture}[x=1.04mm,y=1mm,baseline]
    \node[anchor=west,font=\fontsize{22}{26}\selectfont\bfseries,text=ink] at (0,11) {#1};
    \fill[lightgray] (0,0) rectangle (100,7);
    \fill[#4] (0,0) rectangle (#2,7);
    \node[anchor=west,font=\fontsize{21}{25}\selectfont\bfseries,text=ink] at (102,3.5) {#3};
  \end{tikzpicture}\par\vspace{2.2mm}
}

\begin{document}
\bodyfont

% ============================================================
% HEADER
% ============================================================
\begin{center}
  \setlength{\parskip}{0.8mm}
  \IfFileExists{../assets/university-trier.pdf}{%
    \includegraphics[width=66mm]{../assets/university-trier.pdf}\par
  }{%
    {\fontsize{29}{34}\selectfont\bfseries University of Trier}\par
  }

  \vspace{2mm}
  {\fontsize{44}{49}\selectfont\bfseries Static vs Contextual Embeddings for Lexical Ambiguity\par}
  {\fontsize{29}{34}\selectfont\textcolor{trierblue}{A Word in Context Evaluation}\par}

  \vspace{3mm}
  \begin{minipage}[t]{0.43\textwidth}\centering
    {\fontsize{25}{29}\selectfont\bfseries \FirstAuthor}\par
    {\fontsize{21}{25}\selectfont Enrollment no. \FirstEnrollment}
  \end{minipage}
  \hspace{7mm}
  \begin{minipage}[t]{0.43\textwidth}\centering
    {\fontsize{25}{29}\selectfont\bfseries \SecondAuthor}\par
    {\fontsize{21}{25}\selectfont Enrollment no. \SecondEnrollment}
  \end{minipage}\par

  \vspace{1mm}
  {\fontsize{20.5}{24}\selectfont University of Trier \quad | \quad MSc Natural Language Processing \quad | \quad Trends in Natural Language Processing}\par
  {\fontsize{18}{22}\selectfont\textcolor{trierblue}{github.com/Mr-Dark-debug/static-contextual-ambiguity}}
\end{center}

\vspace{-1mm}
{\color{midgray}\rule{\textwidth}{0.7pt}}
\vspace{4mm}

% ============================================================
% Balanced columns let the explanatory text continue naturally between columns.
% ============================================================
\begin{multicols}{3}
\RaggedRight

\sectiontitle{Introduction}
\keybox{%
``She sat on the \textbf{\textit{bank}} of the river.'' Here, \textit{bank} means the land beside the water.\\[1.5mm]
``She deposited money at the \textbf{\textit{bank}}.'' Here, \textit{bank} means a financial institution.
}
\par\vspace{3mm}

The spelling is identical, but the meaning changes because of the surrounding words. This is \textbf{lexical ambiguity}.

We wanted to understand whether a computer model represents the same word differently when its context changes, and whether that difference helps it decide if the meaning is the same or different.

\sectiontitle{Inspiration}
Our experiment follows the \textbf{Word in Context (WiC)} task introduced by Pilehvar and Camacho-Collados (2019). WiC avoids asking for a dictionary definition. It asks \textbf{whether the same word has the same meaning in two sentences}.

Camacho-Collados and Pilehvar (2018) describe a problem with traditional word representations. One stored vector can mix several meanings together. That motivated our comparison between a fixed representation and representations that use context.

\sectiontitle{Background}
A \textbf{word embedding} is a list of numbers used to represent a word computationally. We compare two pretrained approaches called \textbf{GloVe} and \textbf{BERT}.

\textbf{GloVe} learns one stored vector for each word from patterns of words occurring together in large text collections. The vector for \textit{bank} therefore stays the same in a river sentence and a money sentence. This is a \textbf{static embedding}.

\textbf{BERT} reads the surrounding sentence before creating the target-word representation. Its representation of \textit{bank} can therefore change between those two sentences. This is a \textbf{contextual representation}.

The two meanings of \textit{bank} above are \textbf{homonyms} because they are unrelated. Related meanings are called \textbf{polysemy}. For example, \textit{paper} can mean writing material or a newspaper. The distinction is between unrelated and related meanings.

WiC does not ask the model to name a sense. It only asks whether the target word has the \textbf{same meaning} or a \textbf{different meaning} in the two sentences.

\columnbreak
\sectiontitle{Hypothesis}
We expected more context to make changes in meaning easier to recognise.

GloVe gives \textit{bank} the same vector in both sentences, so it cannot see the difference between the river and money uses. Of the \textbf{638 validation pairs}, 319 have the same meaning and 319 have different meanings. Always choosing one answer would get half correct, so \textbf{50\% is our simple baseline}.

We expected nearby words to help GloVe and the full sentence to help BERT more.

\sectiontitle{Methodology}

\subsectiontitle{The dataset}
We used the \textbf{SuperGLUE Word in Context (WiC)} dataset throughout. Each example has \textbf{two sentences with the same marked word} and a label saying whether its meaning is the same.

A \textit{bank} beside a river paired with a financial \textit{bank} is labelled \textbf{different meaning}. Two financial uses of \textit{bank} would be labelled \textbf{same meaning}.

WiC has \textbf{5,428 labelled training pairs} and \textbf{638 labelled validation pairs}. We chose method settings on training data, then measured performance on validation.

The \textbf{1,400 test pairs} have no public labels, so we report validation scores.

\subsectiontitle{The comparison}
We used \textbf{GloVe and BERT in three conditions}, always on the same data.

\textbf{1. GloVe target word only.} We use its stored vector with 300 numbers. The same word gets the same vector in both sentences.

\textbf{2. GloVe with nearby words.} We average the vectors of up to \textbf{two words before and after the target}, leaving out the target.

\keybox{%
``The \textbf{muddy river} \textit{bank} \textbf{flooded after} rain.''\\[1mm]
For \textit{bank}, the context words are \textbf{muddy, river, flooded, after}.
}
\par\vspace{3mm}

\textbf{3. BERT using the full sentence.} We use a frozen pretrained BERT base model. We combine the target representation from its \textbf{last four layers} without training it further on WiC.

\columnbreak
\subsectiontitle{From vectors to an answer}
For each sentence pair, a method gives us two target representations. We compare them with \textbf{cosine similarity}. A higher score means the representations are more alike. A lower score means they are more different.

Using training data, we select a cutoff for each condition. Scores above the cutoff predict \textbf{same meaning}, while scores below it predict \textbf{different meaning}. Our main score is \textbf{accuracy}, the percentage of the 638 validation pairs answered correctly.

\sectiontitle{Results}
All three conditions were evaluated on the same \textbf{638 validation pairs}.

GloVe using the target word alone answered \textbf{\TargetCorrect{} pairs correctly (\TargetAccuracy)}. Adding nearby GloVe words increased this to \textbf{\ContextCorrect{} pairs (\ContextAccuracy)}. BERT performed best with \textbf{\BertCorrect{} correct pairs (\BertAccuracy)}.

\vspace{2mm}
{\fontsize{27}{32}\selectfont\bfseries\textcolor{deepblue}{Validation accuracy}}\par\vspace{3mm}

\begin{minipage}{\linewidth}
\resultbar{GloVe target only}{\TargetBarValue}{\TargetAccuracy{} \; (\TargetCorrect/638)}{midgray}
\resultbar{GloVe nearby words}{\ContextBarValue}{\ContextAccuracy{} \; (\ContextCorrect/638)}{trierblue!60}
\resultbar{BERT in full sentence}{\BertBarValue}{\BertAccuracy{} \; (\BertCorrect/638)}{trierblue}
\end{minipage}
\par\vspace{2mm}

\keybox{%
{\fontsize{30}{35}\selectfont\bfseries\textcolor{deepblue}{+\AccuracyGain{} percentage points}}\\[1mm]
Compared with GloVe using nearby words, BERT answered \textbf{\BertGainCorrect{} additional sentence pairs correctly}.
}
\par\vspace{3mm}

We repeated that comparison on \textbf{1,000 resampled versions} of the validation set. The estimated 95\% interval for BERT's advantage was \textbf{\GainLower{} to \GainUpper{} percentage points}, so the advantage remained positive throughout this interval.

Looking pair by pair, both stronger methods were correct on \textbf{\BothCorrectCount{} pairs}. BERT alone was correct on \textbf{\BertOnlyCount}, GloVe with nearby words alone on \textbf{\StaticOnlyCount}, and both were wrong on \textbf{\BothWrongCount}. BERT was stronger on this benchmark, but neither method solved every example.

\sectiontitle{Conclusion}
Our results support the idea that \textbf{context helps a model recognise a word's meaning}. The fixed GloVe vector could not distinguish the two uses of a word. Nearby words helped a little, and BERT's representation of the full sentence did better.

BERT still missed \textbf{\BertErrorCount{} of 638 pairs}. We found a clear improvement, not a complete solution. Building and checking this comparison ourselves made the difference between static and contextual embeddings much easier to see.

\end{multicols}

\end{document}
